\section{Homological preliminaries}\label{sec:homological}

The tools in this section serve two passages in the construction.
Complete-resolution comparison will turn the triangular cone's
endomorphism calculation into the vanishing needed for transfer to a
symmetric algebra. Stable duality will supply the negative-degree maps
used in that calculation. Separately, projective models and their stable
images will turn the derived-induction object into a finite bimodule,
preserving its maps and connecting morphisms under evaluation.

Throughout the construction, $k$ has characteristic two. Unless explicitly
stated otherwise, algebras and modules are finite dimensional over $k$.
Resolution and model complexes have finite-dimensional terms; the full Hom
complexes defined below may have infinite-dimensional product terms.
Modules are left modules unless a side is specified. Write
$D=\Hom_k(-,k)$ and $G^e=G\otimes_kG^{\mathrm{op}}$; an unsubscripted tensor
product is over $k$. Differentials lower homological degree, and
\[
 P[a]_j=P_{j-a},\qquad c_jP=\coker(d^P_{j+1}:P_{j+1}\longrightarrow P_j).
\]
Characteristic two removes the signs in shifts, tensor differentials and cones.
For two complexes we use the full Hom complex
\begin{equation}\label{h:full-hom}
 \Hom_G^a(P,Q)=\prod_{j\in\mathbb Z}\Hom_G(P_j,Q_{j-a}),\qquad
 (\delta f)_j=d^Q_{j-a}f_j+f_{j-1}d^P_j.
\end{equation}
For a module $N$, viewed in degree zero, $\Hom_G(P,N)^a=\Hom_G(P_a,N)$.
The space $\sHom_G(M,N)$ is the quotient of $\Hom_G(M,N)$ by maps factoring
through finite projectives.

\subsection{Complete resolutions and stable duality}

An exact complex $P$ of projectives is \emph{totally acyclic} if
$\Hom_G(P,G)$ is exact. If $c_0P=M$, we call $P$ a complete resolution of $M$.
The comparison below also applies when $G$ is not self-injective.
Complete resolutions and passage from derived objects to stable modules
belong to the framework developed by Buchweitz~\cite{Buchweitz2021}.
We give the comparison and truncation arguments in our conventions.

\begin{lemma}[Complete-resolution comparison]\label{h:comparison}
Let $P$ be totally acyclic and let $Q$ be an exact complex of projectives.
For every integer $a$ there are natural isomorphisms
\begin{equation}\label{h:comparison-formula}
 H^a\Hom_G^\bullet(P,Q)
 \simeq\sHom_G(c_aP,c_0Q)
 \simeq H^a\Hom_G(P,c_0Q).
\end{equation}
The first isomorphism is induced by cokernels of shifted chain maps and
respects composition. For $a>0$ these spaces are
$\Ext_G^a(c_0P,c_0Q)$.
\end{lemma}

\begin{proof}
We first record the extension property supplied by total acyclicity. If $J$
is projective, then $\Hom_G(P,J)$ is exact, since $J$ is a summand of a finite
sum of copies of $G$. Identify $c_jP$ with $\im d_j^P\subset P_{j-1}$.
For $u:c_jP\to J$, the map $u d_j^P:P_j\to J$ is a cocycle.
Exactness gives $h:P_{j-1}\to J$ with $h d_j^P=u d_j^P$, so $h$ extends $u$.
It follows at once that the cycles of $\Hom_G(P,N)^a$ are maps
$c_aP\to N$ and its boundaries are exactly the maps factoring through
projectives. This proves the second isomorphism, with $N=c_0Q$.

For the first, begin in degree zero. A map $c_0P\to c_0Q$ lifts to
$f_0:P_0\to Q_0$, and then successively to $f_j$ for $j>0$, using
projectivity of $P_j$ and exactness of $Q$. To proceed downwards,
$d_j^Qf_j$ vanishes on $\ker d_j^P=\im d_{j+1}^P$ by the chain equation
already obtained in degree $j+1$. It therefore defines a map
$c_jP=\im d_j^P\to Q_{j-1}$. The extension property gives $f_{j-1}$.
Thus every cokernel map lifts to a full chain map.

A homotopy induces a projective factorization on $c_0$. Conversely, suppose
the cokernel map of $f$ factors as $c_0P\to J\to c_0Q$ with $J$ projective.
Extend its first factor to $P_{-1}$ and lift its second factor to $Q_0$.
Their composite is a homotopy component $h_{-1}:P_{-1}\to Q_0$ whose
homotopy differential removes the induced map on cokernels. It remains to
nullhomotope a chain map inducing zero there. For this remaining map set
$h_{-1}=0$. For $j\geq0$, lift
\[
 f_j+h_{j-1}d_j^P:P_j\longrightarrow\im d_{j+1}^Q
\]
to $h_j:P_j\to Q_{j+1}$; the initial image assertion is the vanishing on
$c_0$, and each later one follows from the preceding chain equation.
For $j<0$, the map $f_j+d_{j+1}^Qh_j$ kills $\im d_{j+1}^P$ by the equation
in degree $j+1$. It descends to $\im d_j^P$ and extends to
$h_{j-1}:P_{j-1}\to Q_j$ by the same extension property. These equations
give $f=dh+hd$ in every degree. The product in \eqref{h:full-hom} permits all
these components; each equation uses only two of them.

Replace $P$ by $P[-a]$ to obtain degree $a$. The construction by cokernels
preserves compositions and shifts. Finally, the nonnegative tail of $P$
resolves $c_0P$. In every positive cohomological degree its Hom complex has
the same two adjacent differentials as $\Hom_G(P,c_0Q)$, proving the Ext
assertion, including degree one.
\end{proof}

Recall that $G$ is \emph{symmetric} if there is a bimodule isomorphism
$G\simeq DG$, or equivalently a linear form $t:G\to k$ such that
$t(ab)=t(ba)$ and $(a,b)\mapsto t(ab)$ is nondegenerate. In this case finite
projectives are injective, and every finite module embeds in a finite
projective: dualize a finite projective surjection onto its dual.
Splicing a projective resolution and a projective-injective coresolution
therefore gives a complete resolution of every finite module; applying
$\Hom_G(-,G)$ preserves its exactness because $G$ is injective.

For such a resolution put $M[a]=c_{-a}P$ in the stable category. The exact
sequences
\begin{equation}\label{h:syzygy-sequence}
 0\longrightarrow c_{j+1}P\longrightarrow P_j
 \longrightarrow c_jP\longrightarrow0
\end{equation}
show that $[1]$ is inverse syzygy. Lemma~\ref{h:comparison} makes these
shifts independent of choices, functorial on stable maps, and compatible
with composition. In particular,
\begin{equation}\label{h:stable-ext}
 H^a\Hom_G(P,N)=\sHom_G(M,N[a])\qquad(a\in\mathbb Z).
\end{equation}
For positive degrees the shifted composition is the Yoneda product, since
both are computed by lifting maps to projective resolutions.

A short exact sequence $0\to L\to V\to N\to0$ gives a long exact sequence
of the groups in \eqref{h:stable-ext}: apply $\Hom_G(P,-)$ termwise.
Its connecting map is composition with $\partial:N\to L[1]$. Concretely,
lift a projective surjection $J\to N$ to $J\to V$ and restrict to its kernel.
The resulting map $\Omega N\to L$, shifted once, is $\partial$.
This agrees with the cochain connecting map, which lifts a cocycle and
then applies the differential. We use this description to keep track of
connecting maps, as well as the groups in the long exact sequences.

The next lemma is the symmetric form of Auslander--Reiten duality,
also called Tate duality; see
\cite[Section~2, (2.3) and (2.7)--(2.8)]{Linckelmann2013}.
Its composition pairing will determine the action on negative-degree
self-maps.

\begin{lemma}[Symmetric stable duality]\label{h:duality}
For a symmetric algebra $G$ there are natural isomorphisms
\begin{equation}\label{h:duality-formula}
 D\sHom_G(M,N)\simeq\sHom_G(N,M[-1]).
\end{equation}
There is a linear functional $\lambda_M:\sHom_G(M,M[-1])\to k$ for which
the perfect pairing is $(f,g)\mapsto\lambda_M(gf)$. Consequently, for every
integer $a$ the pairing
\begin{equation}\label{h:shifted-pairing}
 \sHom_G(M,N[a])\times\sHom_G(N,M[-a-1])\longrightarrow k,
 \qquad (\alpha,\beta)\longmapsto\lambda_M\bigl(\beta[a]\alpha\bigr)
\end{equation}
is perfect and natural.
\end{lemma}

\begin{proof}
For every finite projective $J$ the finite-dimensional adjunction gives
\[
 D\Hom_G(J,N)\simeq\Hom_G(N,DG\otimes_GJ)\simeq\Hom_G(N,J).
\]
The first identification is immediate for $J=G$, hence for finite sums
and summands; the second uses a fixed symmetrizing form. Apply this
termwise to a complete resolution $P$ of $M$. Dualizing degree-zero
cohomology of $\Hom_G(P,N)$ yields
\[
 \frac{\Hom_G\bigl(N,\ker(P_0\to P_{-1})\bigr)}
 {\im\bigl(\Hom_G(N,P_1)\to
                \Hom_G(N,\ker(P_0\to P_{-1}))\bigr)}.
\]
Here $\ker(P_0\to P_{-1})=c_1P=M[-1]$. A factorization through any
projective lifts through $P_1\twoheadrightarrow c_1P$, so the displayed
quotient is precisely $\sHom_G(N,M[-1])$. All identifications are natural;
comparison of complete resolutions also gives naturality in $M$.

Define $\lambda_M(g)$ by pairing $1_M$ with $g$. Naturality in $N$, applied
to $f:M\to N$, identifies the pairing of $f$ and $g$ with
$\lambda_M(gf)$. Replacing $N$ by $N[a]$ and using the shift equivalences
gives \eqref{h:shifted-pairing}.
\end{proof}

\subsection{Projective models and their stable images}

In Section~\ref{b:section}, derived induction first supplies a bounded
derived bimodule. We will obtain the finite bimodule $Y$ from high
cokernels of a projective model. To compute its evaluated stable Hom
groups, we must also retain the actual maps attaching the model's
successive pieces. The next two lemmas establish this passage from
derived objects and connecting morphisms to finite stable modules.

We write $\mathrm D^b(G)$ for the derived category of bounded complexes of
finite modules. A complex is \emph{perfect} if it is derived isomorphic to
a bounded complex of finite projectives. We will represent bounded derived
objects by bounded-below complexes of finite projectives with bounded
homology. The following elementary facts specify the models and maps we use.
For a chain map $f:L\to N$, our cone convention is
\begin{equation}\label{h:cone}
 \Cone(f)_j=N_j\oplus L_{j-1},\qquad
 d_j^{\Cone(f)}=
 \begin{pmatrix}d_j^N&f_{j-1}\\0&d_{j-1}^L\end{pmatrix}.
\end{equation}

The projective-model statement is the bounded case of the standard
$K$-projective formalism; see
\cite[Proposition~1.4 and Example~3.2(a)]{Spaltenstein1988}, with the
opposite indexing convention for complexes.

\begin{lemma}[Derived maps and two-cell models]\label{h:derived-models}
A bounded-below complex $P$ of projectives computes derived maps out of it
by chain maps modulo homotopy. Every bounded complex of finite modules has
such a projective model with finite terms. In particular, two such models
of the same derived object are homotopy equivalent.

Suppose a derived triangle is given by
\[
 L\longrightarrow V\longrightarrow R\xrightarrow{\partial}L[1].
\]
For projective models $P_L,P_R$ and a chain map
$b:P_R\to P_L[1]$ representing $\partial$, a model of $V$ is the two-cell
complex
\begin{equation}\label{h:two-cell}
 P_{V,j}=P_{L,j}\oplus P_{R,j},\qquad
 d_j^{P_V}=\begin{pmatrix}d_j^{P_L}&b_j\\0&d_j^{P_R}\end{pmatrix}.
\end{equation}
Its subcomplex and quotient represent $L$ and $R$, and its connecting
map is the specified $\partial$.
\end{lemma}

\begin{proof}
Every chain map from $P$ to an acyclic complex is nullhomotopic: start at
the lowest nonzero degree of $P$ and successively lift the residual map
through the next differential of the target. The same argument applies
to every shift. Thus the full Hom from $P$ sends quasi-isomorphisms to
quasi-isomorphisms, which proves the assertion about derived maps. This is
the $K$-projectivity of a bounded-below projective source; it does not assert
$K$-projectivity of a complete resolution.

For completeness, a bounded-below finite complex $V$ admits a projective
model $\theta:Q\to V$ by the following induction. Having constructed lower
terms, choose $Q_j$ finite projective surjecting onto
\[
 \{(x,y)\in V_j\oplus Q_{j-1}:d_Vx=\theta(y),\ d_Qy=0\}.
\]
Use its two components for $\theta_j$ and $d_j^Q$. Every cycle of the
comparison cone is then a boundary, so $\theta$ is a quasi-isomorphism.
A derived isomorphism between projective models and its inverse are
represented by chain maps, and their composites are homotopic to the
identities by the first assertion.

To verify the last assertion, represent the first arrow of the triangle
by $f:P_L\to P_V$ on any projective models. Its cone represents $R$, and
the projection $\Cone(f)\to P_L[1]$ represents $\partial$.
The cone of the deshifted projection has $P_V$ as a subcomplex with
contractible quotient $\Cone(1_{P_L})$. It therefore represents $V$.
Replacing the projective models by homotopy equivalent ones and the map
by its prescribed representative preserves this conclusion: a homotopy
commuting square induces a cone map with that homotopy as its off-diagonal
component, and homotopic maps have isomorphic cones. The resulting cone
$\Cone(b[-1])$ is exactly \eqref{h:two-cell}.
\end{proof}

A bounded-below complex with bounded homology represents a bounded derived
object by good truncation. Also, a bounded complex whose terms have finite
projective dimension is perfect: its finite filtration by terms expresses
it as successive cones of shifts of modules with bounded projective
resolutions. These observations allow us to use the preceding models
without introducing infinite-dimensional terms.

\begin{lemma}[Stable images]\label{h:stable-image}
Let $G$ be symmetric and let $P$ be bounded below with projective terms and
bounded homology. For $n$ sufficiently large the formula
\begin{equation}\label{h:stable-image-formula}
 \langle P\rangle=(c_nP)[n]
\end{equation}
defines a stable module, independent of $n$ up to natural stable
isomorphism. It is functorial on derived maps represented by such models,
commutes with shifts, sends a module to its stable class, and sends an Ext
class to its corresponding stable map.

A degreewise split short exact sequence $0\to P_L\to P_V\to P_R\to0$
of these complexes gives a stable connecting map
$\langle P_R\rangle\to\langle P_L\rangle[1]$ and the associated long exact
stable Hom sequences. With respect to a graded splitting, this connecting
map is the stable image of the cross differential
$b:P_R\to P_L[1]$ in \eqref{h:two-cell}.
\end{lemma}

\begin{proof}
Above the homology support, \eqref{h:syzygy-sequence} is exact for $P$.
It identifies $(c_{n+1}P)[n+1]$ with $(c_nP)[n]$, naturally on chain maps.
A homotopy changes $c_n(f)$ by a map factoring through $P_{n-1}$, so the
construction descends to homotopy classes. Lemma~\ref{h:derived-models}
then makes it functorial on derived maps. The identity
$c_n(P[a])=c_{n-a}P$ gives shift compatibility. A projective resolution of
a module recovers its successive syzygies, proving the assertion for
modules. Restricting a lifted representative of an Ext class to high
degrees gives the same successive syzygies of that map, proving the last
assertion about classes.

Choose $n$ so that all three complexes are exact in degrees at least $n$.
Their cokernels give a short exact sequence
\[
 0\longrightarrow c_nP_L\longrightarrow c_nP_V
 \longrightarrow c_nP_R\longrightarrow0.
\]
For injectivity on the left, suppose
$(x,0)=d_{n+1}^{P_V}(u,v)$. Then $d_{n+1}^{P_R}v=0$, so
$v=d_{n+2}^{P_R}w$ for some $w$. The equation $d_Lb+bd_R=0$ gives
$x=d_{n+1}^{P_L}(u+b_{n+2}w)$, as required. Surjectivity and exactness in
the middle follow by lifting components in the graded splitting.

There is a commutative diagram with exact rows
\[
\begin{tikzcd}[column sep=large]
0\arrow[r]&c_{n+1}P_R\arrow[r]\arrow[d,"\overline b_{n+1}"']
 &P_{R,n}\arrow[r]\arrow[d]
 &c_nP_R\arrow[r]\arrow[d,equal]&0\\
0\arrow[r]&c_nP_L\arrow[r]&c_nP_V\arrow[r]&c_nP_R\arrow[r]&0.
\end{tikzcd}
\]
The middle vertical map sends $z$ to the class of $(0,z)$. On an element
$d_{n+1}^{P_R}z$, its value is the class of $(b_{n+1}z,0)$, since
$(b_{n+1}z,d_{n+1}^{P_R}z)$ is a boundary in $P_V$. Thus the left vertical
map is induced by $b_{n+1}:P_{R,n+1}\to P_{L,n}$.
Shifting the left vertical arrow by $[n+1]$ gives
$\langle P_R\rangle\to\langle P_L\rangle[1]$. By the description after
\eqref{h:stable-ext}, this is the connecting map of the bottom row shifted
by $[n]$, and it is exactly $\langle b\rangle$. This also gives the
asserted long exact sequences.
\end{proof}

The following criterion supplies the other use of stable images in
Section~\ref{b:section}. There, an ordinary trace will rule out a perfect
factorization of a derived map. The criterion then shows that its stable
image is nonzero.

\begin{lemma}[Vanishing stable image and perfect factorization]
\label{h:perfect-factor}
Over a symmetric algebra, a derived map between bounded derived objects
whose stable image is zero factors in the derived category through a
perfect complex.
\end{lemma}

\begin{proof}
Choose bounded-below projective models $P,Q$ and represent the map by a
global chain map $f:P\to Q$, using Lemma~\ref{h:derived-models}. For a
sufficiently high $n$, both complexes are exact in degrees at least $n$
and $c_n(f)$ factors through a finite projective $J$:
\[
 c_nP\xrightarrow{u}J\xrightarrow{v}c_nQ.
\]
Since $J$ is injective, extend $u$ across
$c_nP=\im d_n^P\subset P_{n-1}$. Since $J$ is projective, lift $v$ to $Q_n$.
The composite gives $h_{n-1}:P_{n-1}\to Q_n$. It makes
$f_n+h_{n-1}d_n^P$ zero after projection to $c_nQ$, so this residual map
lifts through $d_{n+1}^Q$ to $h_n:P_n\to Q_{n+1}$. Inductively, for $j>n$
the chain equation and the preceding homotopy equation show that
$f_j+h_{j-1}d_j^P$ lands in $\ker d_j^Q=\im d_{j+1}^Q$. Lift it to $h_j$.
Set $h_j=0$ for $j<n-1$.

The homotopic map $f'=f+dh+hd$ now vanishes in every degree $j\geq n$.
Its chain equation in degree $n$ reads
\begin{equation}\label{h:truncation-boundary}
 f'_{n-1}d_n^P=0.
\end{equation}
Consequently its remaining components define a chain map into the actual
homological subcomplex $Q_{\leq n-1}$, obtained by retaining $Q_j$ for
$j\leq n-1$ and setting the higher terms to zero. Equation
\eqref{h:truncation-boundary} verifies the boundary condition at the cut.
Thus $f'$ factors through $Q_{\leq n-1}$, which is bounded and projective
because $Q$ was bounded below. Composing with the chosen model maps gives
the desired derived factorization. This also covers a map with a module
as target: first lift it globally to a projective resolution of that
module, and then apply this argument.
\end{proof}

\subsection{Bimodule evaluation}

For a symmetric algebra $E$, the tensor product of symmetrizing forms makes
$E^e$ symmetric. We call an $E$-bimodule \emph{side-projective} if it is
projective as a left and as a right $E$-module, and write $LX=L\otimes_EX$.
Side-projectivity does not assert projectivity over $E^e$.
For the derived-induction construction, the passage to a finite stable
bimodule must commute with evaluation at $X$, including the attaching
maps just discussed. The next lemma establishes these compatibilities.

\begin{lemma}[Finite shifts and evaluation]\label{h:evaluation}
Every stable shift of a finite side-projective bimodule has a finite
side-projective representative. Evaluation at a finite $E$-module $X$
sends stable maps between these bimodules to stable $E$-module maps and
has natural identifications
\begin{equation}\label{h:evaluation-shifts}
 (L[a])X\simeq(LX)[a],\qquad L(X[a])\simeq(LX)[a].
\end{equation}
It preserves the stable connecting maps of short exact sequences of
side-projective bimodules, and these identifications are natural in both
$L$ and $X$.

If $P$ is a bounded-below complex of projective bimodules with bounded
homology and all sufficiently high $c_nP$ are side-projective, then
$P\otimes_EX$ has bounded homology and
\begin{equation}\label{h:evaluation-image}
 \langle P\rangle X\simeq\langle P\otimes_EX\rangle,
\end{equation}
naturally on maps and on connecting maps from degreewise split sequences
of such complexes. The hypothesis on the high cokernels holds whenever
$P$, considered separately on either side, is homotopy equivalent to a
bounded complex of projectives on that side.
\end{lemma}

\begin{proof}
Take a projective bimodule surjection onto $L$. Its kernel is
side-projective because the sequence splits on each side. For a positive
shift, use the finite projective-injective embedding
\[
 L\longrightarrow\Hom_k(E^e,L),\qquad
 z\longmapsto(a\longmapsto az),
 \qquad (r\varphi)(a)=\varphi(ar).
\]
Its target is injective by coinduction and projective by symmetry of
$E^e$. The embedding splits on each side because $L$ is injective on
that side, so its cokernel is again side-projective. Iteration supplies
finite representatives for every shift.

An $E^e$-projective evaluates to an $E$-projective: this is immediate for
the free bimodule $E\otimes E$, which evaluates to $E\otimes X$, and
passes to finite sums and summands. Hence evaluation respects projective
factorizations. Every short exact sequence of side-projective bimodules
splits on the right, and therefore remains exact after evaluation. In
particular, evaluating the sequences just used to define shifts gives
the first identification in \eqref{h:evaluation-shifts}. Applying the
projective-presentation description of connecting maps to these same
sequences shows that their actual connecting maps are preserved.

For fixed $L$, right projectivity makes $L\otimes_E-$ exact, and left
projectivity makes it preserve projectives. Applying it to projective
resolutions and projective-injective coresolutions gives the second
identification in \eqref{h:evaluation-shifts}. A bimodule map $L\to L'$
gives a natural transformation on these complexes; thus evaluating it
on a shifted module agrees stably with shifting its evaluation. The
same reasoning proves naturality in $X$ and compatibility with
composition.

For the complex assertion, the high exact sequences
$0\to c_{n+1}P\to P_n\to c_nP\to0$ split on the right. Their evaluations
remain exact, proving bounded homology of $P\otimes_EX$; its terms are
projective by the preceding argument. Right exactness gives the
canonical identifications
\[
 c_n(P\otimes_EX)=(c_nP)\otimes_EX.
\]
Together with \eqref{h:evaluation-shifts}, these identify the two sides
of \eqref{h:evaluation-image}. They commute with chain maps. For a
filtered complex, the cross differential in
Lemma~\ref{h:stable-image} evaluates to its tensor with $X$. That lemma's
explicit diagram therefore identifies the evaluated connecting map with
the connecting map after evaluation.

Finally, suppose a homotopy equivalence on one side compares $P$ with a
complex supported in degrees at most $d$. At any $n>d$, the homotopy
between the identity and the composite through that complex gives
\[
 1_{P_n}=d_{n+1}h_n+h_{n-1}d_n.
\]
Writing $\pi_n:P_n\to c_nP$, passage to the cokernel yields
\[
 1_{c_nP}=(\pi_nh_{n-1})\overline d_n,
 \qquad \overline d_n:c_nP\longrightarrow P_{n-1}.
\]
Thus $c_nP$ is a summand of a projective on that side. The separate
homotopy equivalence on the other side proves the other projectivity.
The homotopies need not be bimodule homotopies.
\end{proof}
